\chapter*{Resumen}
\addcontentsline{toc}{chapter}{Resumen}

Este Trabajo de Fin de Grado presenta el diseño, desarrollo y despliegue de un sistema de Generación Aumentada mediante Recuperación (\textit{Retrieval-Augmented Generation}, RAG) de ámbito corporativo, ejecutado íntegramente en infraestructura local (\textit{on-premise}) e integrando herramientas y modelos de código abierto con desarrollo propio.

El sistema permite consultar en lenguaje natural los documentos internos de una organización, páginas web y fuentes legislativas oficiales, respondiendo siempre con las fuentes citadas. Una pila de microservicios contenerizados con Docker integra modelos de lenguaje locales (Ollama), una base de datos vectorial (Qdrant) y un agente encaminador denominado JARVIS que detecta la intención de cada consulta y la delega al modo de operación más adecuado. Entre las fuentes integradas destacan Microsoft SharePoint ---principal entorno de validación práctica---, páginas web extraídas mediante scraping, documentos escaneados tratados con OCR acelerado por GPU y el Boletín Oficial del Estado.

El sistema ha sido desplegado y validado en un entorno real con resultados satisfactorios: tiempos de respuesta inferiores a 3.5 segundos y una reducción muy significativa de las alucinaciones en modo RAG gracias a una estrategia de recuperación documental orientada a la trazabilidad.

\vspace{0.5cm}
\noindent\textbf{Palabras clave:} RAG, LLM, inteligencia artificial, embeddings, base de datos vectorial, microservicios, Docker, Ollama, Qdrant, on-premise, código abierto.

\newpage

\chapter*{Abstract}
\addcontentsline{toc}{chapter}{Abstract}

This Bachelor's Thesis presents the design, development and deployment of a corporate Retrieval-Augmented Generation (RAG) system, executed entirely on local infrastructure (on-premise) and built by integrating open-source tools and models together with original development.

The system allows the internal documents of an organization, web pages and official legislative sources to be queried in natural language, always answering with cited sources. A stack of Docker-containerized microservices integrates locally executed language models (Ollama), a vector database (Qdrant) and a routing agent called JARVIS, which detects the intent of each query and delegates it to the most suitable operating mode. Integrated sources include Microsoft SharePoint ---the main real-world validation environment---, web pages processed through scraping, scanned documents handled with GPU-accelerated OCR, and the Spanish Official State Gazette (BOE).

The system has been deployed and validated in a real environment with satisfactory results: response times under 3.5 seconds and a very significant reduction of hallucinations in RAG mode thanks to a document retrieval strategy focused on traceability.

\vspace{0.5cm}
\noindent\textbf{Keywords:} RAG, LLM, artificial intelligence, embeddings, vector database, microservices, Docker, Ollama, Qdrant, on-premise, open source.

\newpage
